
\documentclass{article}
\usepackage{colortbl}
\usepackage{makecell}
\usepackage{multirow}
\usepackage{supertabular}

\begin{document}

\newcounter{utterance}

\centering \large Interaction Transcript for game `wordle', experiment `high\_frequency\_words\_no\_clue\_no\_critic', episode 7 with deepseek{-}v4{-}pro.
\vspace{24pt}

{ \footnotesize  \setcounter{utterance}{1}
\setlength{\tabcolsep}{0pt}
\begin{supertabular}{c@{$\;$}|p{.15\linewidth}@{}p{.15\linewidth}p{.15\linewidth}p{.15\linewidth}p{.15\linewidth}p{.15\linewidth}}
    \# & \multicolumn{2}{c}{Player} && \multicolumn{2}{c}{Game Master} \\
    \hline

    \theutterance \stepcounter{utterance}  
    & & & \multicolumn{4}{p{0.6\linewidth}}{
        \cellcolor[rgb]{0.9,0.9,0.9}{
            \makecell[{{p{\linewidth}}}]{
                \texttt{\tiny{[P1$\langle$GM]}}
                \texttt{Eres un mago del lenguaje al que le gusta adivinar palabras usando las reglas dadas.} \\
\\ 
\texttt{¡Bienvenido a Wordle! Tienes seis intentos para adivinar la palabra objetivo, una palabra válida en español de cinco letras minúsculas. Por favor, usa las etiquetas "explicación:" e "intento:" para proporcionar una explicación concisa para cada intento.} \\
\\ 
\texttt{Por ejemplo, si tu intento es "parte", tu respuesta debería ser} \\
\texttt{explicación: esta es una palabra común en español de cinco letras, y estoy comenzando mi intento con esta palabra.} \\
\texttt{intento: parte} \\
\\ 
\texttt{Después de cada intento, tu respuesta será validada y recibirás retroalimentación indicando qué letras son correctas (verde), qué letras son correctas pero están en la posición incorrecta (amarillo) y qué letras son incorrectas (rojo). Esta retroalimentación puede ser útil para determinar qué letras incluir o excluir en tu siguiente intento.} \\
\\ 
\texttt{Por ejemplo, la retroalimentación para "parte" podría ser:} \\
\texttt{retroalimentación\_intento: p<amarillo> a<amarillo> r<verde> t<amarillo> e<rojo>} \\
\\ 
\texttt{La explicación debe contener detalles sobre cómo se usa el retroalimentación\_intento para llegar a un nuevo intento.} \\
\\ 
\texttt{No generes ningún otro texto.} \\
\\ 
\texttt{Comencemos con tu primer intento.} \\
            }
        }
    }
     \\ \\

    \theutterance \stepcounter{utterance}  
    & \multicolumn{4}{p{0.6\linewidth}}{
        \cellcolor[rgb]{0.561,0.737,0.561}{
            \makecell[{{p{\linewidth}}}]{
                \texttt{\tiny{[P1$\rangle$GM]}}
                \texttt{explicación: comenzaré con una palabra común en español de cinco letras que tiene vocales y consonantes frecuentes para obtener buena información inicial.  } \\
\texttt{intento: antes} \\
            }
        }
    }
    & & \\ \\

    \theutterance \stepcounter{utterance}  
    & & & \multicolumn{2}{p{0.3\linewidth}}{
        \cellcolor[rgb]{0.9,0.9,0.9}{
            \makecell[{{p{\linewidth}}}]{
                \texttt{\tiny{[GM$|$GM]}}
                \texttt{attempts = 1} \\
\texttt{target\_word = serie} \\
\texttt{guess = antes} \\
\texttt{guess\_feedback = a<red> n<red> t<red> e<yellow> s<yellow>} \\
            }
        }
    }
    & & \\ \\

    \theutterance \stepcounter{utterance}  
    & & & \multicolumn{4}{p{0.6\linewidth}}{
        \cellcolor[rgb]{0.9,0.9,0.9}{
            \makecell[{{p{\linewidth}}}]{
                \texttt{\tiny{[P1$\langle$GM]}}
                \texttt{retroalimentación\_intento: a<rojo> n<rojo> t<rojo> e<amarillo> s<amarillo>} \\
\\ 
\texttt{Proporciona tu respuesta solo en este formato.} \\
\texttt{explicación: detalles} \\
\texttt{intento: palabra} \\
            }
        }
    }
     \\ \\

    \theutterance \stepcounter{utterance}  
    & \multicolumn{4}{p{0.6\linewidth}}{
        \cellcolor[rgb]{0.561,0.737,0.561}{
            \makecell[{{p{\linewidth}}}]{
                \texttt{\tiny{[P1$\rangle$GM]}}
                \texttt{explicación: la retroalimentación muestra que "a", "n" y "t" no están en la palabra (rojo), mientras que "e" y "s" están presentes pero en posiciones incorrectas (amarillo). Debo encontrar una palabra de cinco letras que contenga "e" y "s" en posiciones diferentes a la cuarta y quinta respectivamente, y sin "a", "n", "t". Probaré con una palabra que tenga "e" y "s" en otras posiciones.} \\
\texttt{intento: verse} \\
            }
        }
    }
    & & \\ \\

    \theutterance \stepcounter{utterance}  
    & & & \multicolumn{2}{p{0.3\linewidth}}{
        \cellcolor[rgb]{0.9,0.9,0.9}{
            \makecell[{{p{\linewidth}}}]{
                \texttt{\tiny{[GM$|$GM]}}
                \texttt{attempts = 2} \\
\texttt{target\_word = serie} \\
\texttt{guess = verse} \\
\texttt{guess\_feedback = v<red> e<green> r<green> s<yellow> e<green>} \\
            }
        }
    }
    & & \\ \\

    \theutterance \stepcounter{utterance}  
    & & & \multicolumn{4}{p{0.6\linewidth}}{
        \cellcolor[rgb]{0.9,0.9,0.9}{
            \makecell[{{p{\linewidth}}}]{
                \texttt{\tiny{[P1$\langle$GM]}}
                \texttt{retroalimentación\_intento: v<rojo> e<verde> r<verde> s<amarillo> e<verde>} \\
\\ 
\texttt{Proporciona tu respuesta solo en este formato.} \\
\texttt{explicación: detalles} \\
\texttt{intento: palabra} \\
            }
        }
    }
     \\ \\

    \theutterance \stepcounter{utterance}  
    & \multicolumn{4}{p{0.6\linewidth}}{
        \cellcolor[rgb]{0.561,0.737,0.561}{
            \makecell[{{p{\linewidth}}}]{
                \texttt{\tiny{[P1$\rangle$GM]}}
                \texttt{explicación: la retroalimentación muestra que "e" en segunda posición, "r" en tercera y "e" en quinta son correctas (verde). La "s" es correcta pero está en posición incorrecta (amarillo), y no puede estar en cuarta posición. La "v" es incorrecta (rojo). La palabra tiene la forma \_ e r \_ e, contiene "s" en primera o segunda posición (pero segunda ya está ocupada por "e"), así que "s" debe ir en primera posición: s e r \_ e. La letra faltante no puede ser "a", "n", "t", "v". Podría ser "s e r v e" pero "v" es incorrecta. Pruebo con "s e r p e".} \\
\texttt{intento: serpe} \\
            }
        }
    }
    & & \\ \\

    \theutterance \stepcounter{utterance}  
    & & & \multicolumn{2}{p{0.3\linewidth}}{
        \cellcolor[rgb]{0.9,0.9,0.9}{
            \makecell[{{p{\linewidth}}}]{
                \texttt{\tiny{[GM$|$GM]}}
                \texttt{Error: The guessed word is not a valid word for this game.} \\
            }
        }
    }
    & & \\ \\

    \theutterance \stepcounter{utterance}  
    & & & \multicolumn{4}{p{0.6\linewidth}}{
        \cellcolor[rgb]{0.9,0.9,0.9}{
            \makecell[{{p{\linewidth}}}]{
                \texttt{\tiny{[P1$\langle$GM]}}
                \texttt{Tu intento no es una palabra válida para este juego. Por favor, inténtalo de nuevo.} \\
\\ 
\texttt{Proporciona tu respuesta solo en este formato.} \\
\texttt{explicación: detalles} \\
\texttt{intento: palabra} \\
            }
        }
    }
     \\ \\

    \theutterance \stepcounter{utterance}  
    & \multicolumn{4}{p{0.6\linewidth}}{
        \cellcolor[rgb]{0.561,0.737,0.561}{
            \makecell[{{p{\linewidth}}}]{
                \texttt{\tiny{[P1$\rangle$GM]}}
                \texttt{explicación: "serpe" no es válida. La palabra tiene la forma s e r \_ e, y la letra faltante no puede ser "a", "n", "t", "v". Probando con una consonante común que forme una palabra válida. "s e r b e" podría funcionar, ya que "serbe" es una conjugación del verbo serbar o podría ser válida. Otra opción es "s e r c e". Probaré con "s e r v e" aunque "v" fue rojo antes, pero tal vez mi último intento no se registró correctamente, lo descarto. Mejor pruebo "s e r c e".} \\
\texttt{intento: serce} \\
            }
        }
    }
    & & \\ \\

    \theutterance \stepcounter{utterance}  
    & & & \multicolumn{2}{p{0.3\linewidth}}{
        \cellcolor[rgb]{0.9,0.9,0.9}{
            \makecell[{{p{\linewidth}}}]{
                \texttt{\tiny{[GM$|$GM]}}
                \texttt{Error: The guessed word is not a valid word for this game.} \\
            }
        }
    }
    & & \\ \\

    \theutterance \stepcounter{utterance}  
    & & & \multicolumn{4}{p{0.6\linewidth}}{
        \cellcolor[rgb]{0.9,0.9,0.9}{
            \makecell[{{p{\linewidth}}}]{
                \texttt{\tiny{[P1$\langle$GM]}}
                \texttt{Tu intento no es una palabra válida para este juego. Por favor, inténtalo de nuevo.} \\
\\ 
\texttt{Proporciona tu respuesta solo en este formato.} \\
\texttt{explicación: detalles} \\
\texttt{intento: palabra} \\
            }
        }
    }
     \\ \\

    \theutterance \stepcounter{utterance}  
    & \multicolumn{4}{p{0.6\linewidth}}{
        \cellcolor[rgb]{0.561,0.737,0.561}{
            \makecell[{{p{\linewidth}}}]{
                \texttt{\tiny{[P1$\rangle$GM]}}
                \texttt{explicación: "serce" no es válida. Con la estructura s e r \_ e, necesito una cuarta letra que forme una palabra en español de cinco letras. Descartadas "a", "n", "t", "v" por rojo, y "p", "c" no formaron palabra. Pruebo "s e r r e", cerrando con "r".} \\
\texttt{intento: serre} \\
            }
        }
    }
    & & \\ \\

    \theutterance \stepcounter{utterance}  
    & & & \multicolumn{2}{p{0.3\linewidth}}{
        \cellcolor[rgb]{0.9,0.9,0.9}{
            \makecell[{{p{\linewidth}}}]{
                \texttt{\tiny{[GM$|$GM]}}
                \texttt{Error: The guessed word is not a valid word for this game.} \\
            }
        }
    }
    & & \\ \\

    \theutterance \stepcounter{utterance}  
    & & & \multicolumn{4}{p{0.6\linewidth}}{
        \cellcolor[rgb]{0.9,0.9,0.9}{
            \makecell[{{p{\linewidth}}}]{
                \texttt{\tiny{[P1$\langle$GM]}}
                \texttt{Tu intento no es una palabra válida para este juego. Por favor, inténtalo de nuevo.} \\
\\ 
\texttt{Proporciona tu respuesta solo en este formato.} \\
\texttt{explicación: detalles} \\
\texttt{intento: palabra} \\
            }
        }
    }
     \\ \\

    \theutterance \stepcounter{utterance}  
    & \multicolumn{4}{p{0.6\linewidth}}{
        \cellcolor[rgb]{0.561,0.737,0.561}{
            \makecell[{{p{\linewidth}}}]{
                \texttt{\tiny{[P1$\rangle$GM]}}
                \texttt{explicación: "serre" no es válida. Con la estructura s e r \_ e, la única letra que parece formar una palabra común es "s e r u e". Probaré "serue" (del verbo seruar, aunque no es común, podría ser válida en este juego).} \\
\texttt{intento: serue} \\
            }
        }
    }
    & & \\ \\

    \theutterance \stepcounter{utterance}  
    & & & \multicolumn{2}{p{0.3\linewidth}}{
        \cellcolor[rgb]{0.9,0.9,0.9}{
            \makecell[{{p{\linewidth}}}]{
                \texttt{\tiny{[GM$|$GM]}}
                \texttt{Error: The guessed word is not a valid word for this game.} \\
            }
        }
    }
    & & \\ \\

    \theutterance \stepcounter{utterance}  
    & & & \multicolumn{4}{p{0.6\linewidth}}{
        \cellcolor[rgb]{0.9,0.9,0.9}{
            \makecell[{{p{\linewidth}}}]{
                \texttt{\tiny{[P1$\langle$GM]}}
                \texttt{Tu intento no es una palabra válida para este juego. Por favor, inténtalo de nuevo.} \\
\\ 
\texttt{Proporciona tu respuesta solo en este formato.} \\
\texttt{explicación: detalles} \\
\texttt{intento: palabra} \\
            }
        }
    }
     \\ \\

    \theutterance \stepcounter{utterance}  
    & \multicolumn{4}{p{0.6\linewidth}}{
        \cellcolor[rgb]{0.561,0.737,0.561}{
            \makecell[{{p{\linewidth}}}]{
                \texttt{\tiny{[P1$\rangle$GM]}}
                \texttt{explicación: "serue" no es válida. Revisando la retroalimentación anterior: en "verse", "e" (pos2) verde, "r" (pos3) verde, "e" (pos5) verde, "s" amarillo. La "s" no está en cuarta posición, y la primera posición está libre. La estructura correcta es s e r \_ e. Letras descartadas: a, n, t, v. Letras no probadas en cuarta posición: b, c, d, f, g, h, j, k, l, m, ñ, p, q, r, s, t, u, v, w, x, y, z. Ya probé p, c, r, u sin éxito. Podría ser "s e r b e" (del verbo serbar) o "s e r l e" o "s e r m e". Probaré "s e r l e".} \\
\texttt{intento: serle} \\
            }
        }
    }
    & & \\ \\

    \theutterance \stepcounter{utterance}  
    & & & \multicolumn{2}{p{0.3\linewidth}}{
        \cellcolor[rgb]{0.9,0.9,0.9}{
            \makecell[{{p{\linewidth}}}]{
                \texttt{\tiny{[GM$|$GM]}}
                \texttt{attempts = 3} \\
\texttt{target\_word = serie} \\
\texttt{guess = serle} \\
\texttt{guess\_feedback = s<green> e<green> r<green> l<red> e<green>} \\
            }
        }
    }
    & & \\ \\

    \theutterance \stepcounter{utterance}  
    & & & \multicolumn{4}{p{0.6\linewidth}}{
        \cellcolor[rgb]{0.9,0.9,0.9}{
            \makecell[{{p{\linewidth}}}]{
                \texttt{\tiny{[P1$\langle$GM]}}
                \texttt{retroalimentación\_intento: s<verde> e<verde> r<verde> l<rojo> e<verde>} \\
\\ 
\texttt{Proporciona tu respuesta solo en este formato.} \\
\texttt{explicación: detalles} \\
\texttt{intento: palabra} \\
            }
        }
    }
     \\ \\

    \theutterance \stepcounter{utterance}  
    & \multicolumn{4}{p{0.6\linewidth}}{
        \cellcolor[rgb]{0.561,0.737,0.561}{
            \makecell[{{p{\linewidth}}}]{
                \texttt{\tiny{[P1$\rangle$GM]}}
                \texttt{explicación: "serle" muestra "l" en rojo, así que la cuarta letra no es "l". Descartadas: a, n, t, v, p, c, r, u, l. Con estructura s e r \_ e, pruebo con "m": "serme" (infinitivo + pronombre enclítico "me", válido como "ser" + "me" aunque es palabra válida en este juego).} \\
\texttt{intento: serme} \\
            }
        }
    }
    & & \\ \\

    \theutterance \stepcounter{utterance}  
    & & & \multicolumn{2}{p{0.3\linewidth}}{
        \cellcolor[rgb]{0.9,0.9,0.9}{
            \makecell[{{p{\linewidth}}}]{
                \texttt{\tiny{[GM$|$GM]}}
                \texttt{attempts = 4} \\
\texttt{target\_word = serie} \\
\texttt{guess = serme} \\
\texttt{guess\_feedback = s<green> e<green> r<green> m<red> e<green>} \\
            }
        }
    }
    & & \\ \\

    \theutterance \stepcounter{utterance}  
    & & & \multicolumn{4}{p{0.6\linewidth}}{
        \cellcolor[rgb]{0.9,0.9,0.9}{
            \makecell[{{p{\linewidth}}}]{
                \texttt{\tiny{[P1$\langle$GM]}}
                \texttt{retroalimentación\_intento: s<verde> e<verde> r<verde> m<rojo> e<verde>} \\
\\ 
\texttt{Proporciona tu respuesta solo en este formato.} \\
\texttt{explicación: detalles} \\
\texttt{intento: palabra} \\
            }
        }
    }
     \\ \\

    \theutterance \stepcounter{utterance}  
    & \multicolumn{4}{p{0.6\linewidth}}{
        \cellcolor[rgb]{0.561,0.737,0.561}{
            \makecell[{{p{\linewidth}}}]{
                \texttt{\tiny{[P1$\rangle$GM]}}
                \texttt{explicación: "m" también es incorrecta. Descartadas: a, n, t, v, p, c, r, u, l, m. La palabra es s e r \_ e. Pruebo con "s e r b e" (de serbar, aunque no es muy común, podría ser válida).} \\
\texttt{intento: serbe} \\
            }
        }
    }
    & & \\ \\

    \theutterance \stepcounter{utterance}  
    & & & \multicolumn{2}{p{0.3\linewidth}}{
        \cellcolor[rgb]{0.9,0.9,0.9}{
            \makecell[{{p{\linewidth}}}]{
                \texttt{\tiny{[GM$|$GM]}}
                \texttt{Error: The guessed word is not a valid word for this game.} \\
            }
        }
    }
    & & \\ \\

    \theutterance \stepcounter{utterance}  
    & & & \multicolumn{4}{p{0.6\linewidth}}{
        \cellcolor[rgb]{0.9,0.9,0.9}{
            \makecell[{{p{\linewidth}}}]{
                \texttt{\tiny{[P1$\langle$GM]}}
                \texttt{Tu intento no es una palabra válida para este juego. Por favor, inténtalo de nuevo.} \\
\\ 
\texttt{Proporciona tu respuesta solo en este formato.} \\
\texttt{explicación: detalles} \\
\texttt{intento: palabra} \\
            }
        }
    }
     \\ \\

    \theutterance \stepcounter{utterance}  
    & \multicolumn{4}{p{0.6\linewidth}}{
        \cellcolor[rgb]{0.561,0.737,0.561}{
            \makecell[{{p{\linewidth}}}]{
                \texttt{\tiny{[P1$\rangle$GM]}}
                \texttt{explicación: "serbe" no es válida. Letras descartadas: a, n, t, v, p, c, r, u, l, m, b. Quedan d, f, g, h, j, k, ñ, q, s, w, x, y, z. La "s" no puede ir en cuarta posición porque ya está en primera. Pruebo con "s e r d e" (serde no sé si es válida, pero es un apellido o podría aceptarse).} \\
            }
        }
    }
    & & \\ \\

    \theutterance \stepcounter{utterance}  
    & & & \multicolumn{2}{p{0.3\linewidth}}{
        \cellcolor[rgb]{0.9,0.9,0.9}{
            \makecell[{{p{\linewidth}}}]{
                \texttt{\tiny{[GM$|$GM]}}
                \texttt{Error: The response should contain the 'intento:' keyword exactly once.} \\
            }
        }
    }
    & & \\ \\

    \theutterance \stepcounter{utterance}  
    & & & \multicolumn{2}{p{0.3\linewidth}}{
        \cellcolor[rgb]{0.9,0.9,0.9}{
            \makecell[{{p{\linewidth}}}]{
                \texttt{\tiny{[GM$|$GM]}}
                \texttt{game\_result = ABORT} \\
            }
        }
    }
    & & \\ \\

\end{supertabular}
}

\end{document}
